\chapter{总结与展望}
本文论述了一种同时面向工程化程序开发与形式化验证的语言与编译工具链原型，
首先介绍了函数式语言、定理证明器的广泛运用于研究现状，指出程序语言、定理证明教育亟待更新。
为了提供一种教学、实验与工程验证的良好平台，本文正文部分依序按 \Tigris 语言报告、类型系统形式化表述、
后端设计与实现、依赖类型编程实例、\Tigrisd 证明器设计与实现五大部分展开研究，
获得了可运行的编译器产出与可用的证明重写器，提出了 TinyTheorem 解决方案。

在程序语言实践层面，本文继承并扩展 Hindley–Milner 类型系统的工程优势，
通过在工程化实践与更强语言表达力之间进行取舍，
创新性地设计出基于 Skolemization 精细控制的 Higher-ranked types、蕴含的谓词约束、严格 HKT 构造等语言特性，
系统化提升可读、可维护的中型编译器构建能力，增加语言表达力。
在形式化验证侧，本文以精简版的依赖类型 \lambdacalc 为计算模型设计核心代数，采用双向（bidirectional）类型检查，
在单一非直谓类型宇宙下的内核中实现了依赖类型、序数族（归纳类型）、Discriminant Refinement 等特性，
并以存在量词的编码与相关定理的证明展示了本系统的使用方式。

两个子系统一体化但解耦，便于教学课程落地，以下就此提出可能的应用场景。

函数式编程范式教学：使用本项目的编程子系统开展本科生课程，介绍函数式编程实践——相比于使用一个工业语言，
子系统因其工程上的相对简单性，适宜教师在备课时穿插函数式语言的设计与实现，有利于加深学生对于其理解。这一方面，
也与工作面试喜欢考察语言“八股”（例如 JAVA/C++ 内部实现上的知识）不谋而合，为学生做好准备。在此基础上，
可以开设研究生课程，或高年级本科生课程，用于讲解编程语言原理、理论，并可作为课程期末作业（项目）的一种范例实现。这一课程常在欧美大学中称为
\emph{Principals of Programming Languages}.

定理证明器使用与实现：使用本项目的定理证明子系统向数学系、理论计算机方向开展研究生课程，
以 \Tigrisd 入门定理证明，讲解依赖类型定理证明器的数学原理与思考方式，
同时，针对计算机科学的学生选择性的论述其实现细节（例如 WHNF），随着课程深入再切换到主流定理证明器，
例如 Coq（Rocq）、Agda 等。

本研究课题目前的成果仍然存在诸多不足。在正文中，我们已经提到过 \Lambdair 的设计失误，除此之外，
再举例而言，函数式语言、定理证明器因其类型系统的复杂性，
需要配套的编辑环境以减少程序员的使用负担。本课题亟待一个可用的 LSP（Language Server Protocol）实现，
与对应的编辑器插件，提供 intellisense、hovering tooltips 等功能，方便用户使用。这一点在定理证明器尤其重要，
因为其注重交互式，在这方面，Lean、Agda、Coq 的编辑器插件交互性就极强，但实现亦复杂。
而一些可选的语言功能，譬如 GADT（序数族）、完整的 HKT、Arbitrary-ranked types 则有利于编程子系统更加贴近学术研究；
而实现一个完整的 CIC 则有利于定理证明器贴近主流实现的理论基础。这些提示有待进一步研究。